\section*{Introduction}
\addcontentsline{toc}{section}{Introduction}

Almost all students studying in a STEM discipline are at some point taught fundamental techniques of linear algebra.
Even within mathematics curricula these introductions to linear algebra terminate at most around the Jordan Canonical Form. This unfortunately fails to make the enthusiastic students aware of the much larger world of intermediate to advanced linear algebra.

The Weierstrass Canonical Form (WCF) can present a stable stepping stone into the world of further linear algebra.
The reason for this twofold. Firstly it is of interest to those studying applied mathematics, engineering and systems theory due to its practical applications in modeling linear systems.
Secondly it should be of interest to students of pure mathematics thanks to its moderately deep geometric foundations and the algebraic techniques used in its proof.

This article, inspired majorly by \cite{berger2012quasi}, hopes to pick up where elementary knowledge of linear algebra left off - provide an intuitive form for Jordan Canonical form and provide an exposition of the WCF.


\subsection*{Disclaimer}
All proofs presented herein, except when explicitly stated otherwise, are the authors creation and the result of a decent level of trial, tribulation and elation.
Any similarity to existing proofs arise, by perhaps the most beautiful property of mathematics, from the structure of the objects being studied.
However the question framework was provided by Stephan Trenn and help with completing the text in a timely manner and in a presentable form from Qianyi.
Naturally praise should be directed to them and complaints to \href{mailto:lingchenyu2288@gmail.com}{\texttt{lingchenyu2288@gmail.com}}.